\documentclass[12pt]{article}
\usepackage[margin=0.6in,top=1.15in,bottom=0.55in]{geometry}
\usepackage{lmodern}
\usepackage{microtype}
\usepackage{titlesec}
\usepackage{enumitem}
\usepackage[hidelinks]{hyperref}
\usepackage{../../latex/grader-worksheet}

\setlength{\parindent}{0pt}
\setlength{\parskip}{0.25em}
\titleformat{\section}{\large\bfseries\scshape}{}{0pt}{}

\GraderSetup{assignment-id=U1L7LE,template-revision=1}

\begin{document}

\begin{center}
{\LARGE\bfseries Week 7 Lit Example Worksheet}\par\vspace{0.05em}
{\large\scshape Macbeth: Appearance of Proof and Pressure}\par\vspace{0.12em}
\rule{0.78\textwidth}{0.5pt}
\end{center}

\textbf{Purpose:} Apply Week 7's Whately method to \textit{Macbeth}. Diagnose arguments that look stronger than they are.

\textbf{Directions:} Use the Lit Example Reader. Do not retell the plot. Separate emotional force from logical force.

\section*{Part 1: Lady Macbeth's Pressure}

\textbf{Question 1.1}\\[-0.15em]
What conclusion is Lady Macbeth pressing on Macbeth in Passage 1? State it in one clear sentence.

\GraderAnswerBox{Part 1 Q1}{2.15in}

\textbf{Question 1.2}\\[-0.15em]
What gives her speech its appearance of proof? Name at least two pressure moves in the language.

\GraderAnswerBox{Part 1 Q2}{2.35in}

\newpage

\section*{Part 2: Defect in the Reasoning}

\textbf{Question 2.1}\\[-0.15em]
Is the main failure in Lady Macbeth's pressure a false claim, a bad inference, or both? Explain.

\GraderAnswerBox{Part 2 Q1}{2.45in}

\textbf{Question 2.2}\\[-0.15em]
Explain the difference between emotional force and logical force using this exchange.

\GraderAnswerBox{Part 2 Q2}{2.35in}

\newpage

\section*{Part 3: The Plan and the Settlement}

\textbf{Question 3.1}\\[-0.15em]
In Passage 2, what does Lady Macbeth's cover-up plan actually support, and what does it fail to prove?

\GraderAnswerBox{Part 3 Q1}{2.45in}

\textbf{Question 3.2}\\[-0.15em]
When Macbeth says ``I am settled,'' has he gained proof that the murder is right, confidence that they can hide it, or both? Explain.

\GraderAnswerBox{Part 3 Q2}{2.45in}

\newpage

\section*{Part 4: Fallacy Autopsy}

\textbf{Question 4.1}\\[-0.15em]
Perform a full fallacy autopsy on Lady Macbeth's persuasion of Macbeth: (1) conclusion, (2) apparent proof, (3) defect, (4) repaired or narrowed claim that would be fairer.

\GraderAnswerBox{Part 4 Q1}{2.65in}

\textbf{Question 4.2}\\[-0.15em]
In one paragraph, explain why a speech can feel conclusive without proving its point, using these Macbeth passages.

\GraderAnswerBox{Part 4 Q2}{2.45in}

\end{document}
